% ===================================================================
%  Dodatek J2: Formalizacja Ścieżki 9 — self-consistent φ-fixed point
%  TGP v1 — Sesja v41 (2026-03-30)
% ===================================================================
\section{Formalizacja \'Scie\.zki~9: self-consistent $\varphi$-fixed point}%
\label{app:J2-formalizacja}

Niniejszy dodatek zamyka g\l{}\'owny wynik Dodatku~\ref{app:ogon-masy}:
\textbf{self-consistent $\varphi$-fixed point}, w~kt\'orym
z\l{}ota proporcja $\varphi = (1+\sqrt{5})/2$ nie jest
za\l{}o\.zeniem, lecz \textbf{wynika z~r\'ownania ODE solitonu
i~jednego warunku wierno\'sciowego}: reprodukcja $r_{21} = m_\mu/m_e$.

Wyniki opieraj\k{a} si\k{e} na skrypcie \texttt{ex106} (14/14 test\'ow
pozytywnych), kt\'ory zunifikowa\l{} i~rozszerzy\l{} analizy
z~\texttt{ex55}--\texttt{ex66}.

\medskip
\noindent\fbox{\parbox{\linewidth}{%
\textbf{Nota kanoniczna (2026-04-07).}
Niniejszy dodatek u\.zywa Formulacji~B ($V = g^3/3-g^4/4$, $g_0^* = 1{,}24915$).
W~sformu\l{}owaniu kanonicznym: $K(g) = g^4$, $P(g) = (\beta/7)g^7 - (\gamma/8)g^8$,
$g_0^e = 0{,}86941$, brak punktu duchowego.
Algebra Koide'a i~r\'ownania master F1--F12 s\k{a} niezale\.zne od tego wyboru.}}
\medskip

% -------------------------------------------------------------------
\subsection{Self-consistent $\varphi$-fixed point}\label{app:J2-fp}
% -------------------------------------------------------------------

\begin{definition}[Self-consistent $\varphi$-fixed point]\label{def:J2-phi-FP}
Niech $A_{\rm tail}(g_0)$ b\k{e}dzie amplitud\k{a} ogona solitonu
(r\'ownanie~\eqref{eq:J-Atail}) jako funkcj\k{a} parametru strzału
$g_0$. \textbf{$\varphi$-fixed point} to warto\'s\'c $g_0^*$ taka, \.ze:
\begin{equation}\label{eq:J2-FP-def}
  \boxed{\left(\frac{A_{\rm tail}(\varphi \cdot g_0^*)}{A_{\rm tail}(g_0^*)}\right)^{\!4}
  \;=\; r_{21}^{\rm PDG} \;=\; 206{,}768.}
\end{equation}
Fizycznie: je\'sli elektron odpowiada $g_0^e = g_0^*$, a~mion
$g_0^\mu = \varphi \cdot g_0^*$ (hipoteza~\ref{hyp:J-golden-topo}),
to warunek~\eqref{eq:J2-FP-def} wymaga, by stosunek mas
$m_\mu / m_e = (A_\mu / A_e)^4$ reprodukowa\l{}
dok\l{}adnie warto\'s\'c eksperymentaln\k{a}.
\end{definition}

\begin{theorem}[Istnienie i~jednoznaczno\'s\'c $\varphi$-FP]\label{thm:J2-FP}
Niech $R(g_0) = (A_{\rm tail}(\varphi g_0)/A_{\rm tail}(g_0))^4 - r_{21}^{\rm PDG}$.
W\'owczas:
\begin{enumerate}[nosep]
  \item $R(g_0)$ zmienia znak w~przedziale $g_0 \in [1{,}15;\, 1{,}35]$:
    $R(1{,}15) \approx +659$, $R(1{,}35) \approx -119$.
  \item Istnieje dok\l{}adnie jedno $g_0^* \in (1{,}15;\, 1{,}35)$
    spe\l{}niaj\k{a}ce $R(g_0^*) = 0$.
  \item Warto\'s\'c numeryczna:
    \begin{equation}\label{eq:J2-g0star}
      g_0^* = 1{,}24915 \pm 0{,}00001,
      \qquad
      g_0^\mu = \varphi \cdot g_0^* = 2{,}02117.
    \end{equation}
  \item Wynikowy stosunek mas:
    \begin{equation}\label{eq:J2-r21-result}
      \boxed{r_{21} = \left(\frac{A_{\rm tail}(2{,}02117)}{A_{\rm tail}(1{,}24915)}\right)^{\!4}
      = 206{,}77
      \qquad (\text{odch.\ } 0{,}0001\% \text{ od PDG}).}
    \end{equation}
\end{enumerate}
\textbf{Dow\'od}: numeryczny (algorytm Brenta, \texttt{ex106},
tolerancja $10^{-5}$). Ci\k{a}g\l{}o\'s\'c $R(g_0)$ wynika z~ci\k{a}g\l{}ej
zale\.zno\'sci rozwi\k{a}zania ODE od warunk\'ow pocz\k{a}tkowych
(twierdzenie o~ci\k{a}g\l{}ej zale\.zno\'sci od parametr\'ow);
zmiana znaku gwarantuje istnienie zera (tw.\ Bolzana).
Jednoznaczno\'s\'c w~$(1{,}15;\, 1{,}35)$ potwierdzona numerycznie
(brak dodatkowych zer w~skanach $N = 60$ punkt\'ow).

\textbf{Status}: \statuslabel{Twierdzenie} (numeryczne,
weryfikacja \texttt{ex106}, test T4: 14/14 PASS).
\end{theorem}

\begin{remark}[Wsp\'olrzystanie z~wcze\'sniejszymi wynikami]\label{rem:J2-comparison}
Wynik~\eqref{eq:J2-r21-result} poprawia dok\l{}adno\'s\'c
z~1{,}6\% (Dodatek~\ref{app:ogon-masy}, uwaga~\ref{rem:J-kluczowy},
$g_0^e = 1{,}24$, $g_0^\mu = 2{,}00$)
do~$0{,}0001\%$.

Zmiana: wcze\'sniejszy wynik u\.zywa\l{} sztywnych
$g_0^e$, $g_0^\mu$; nowy wynik z~$\varphi$-FP
traktuje $g_0^*$ jako wynik dynamiczny,
okre\'slony warunkiem wierno\'sciowym~\eqref{eq:J2-FP-def}.
Numeryczne $g_0^* = 1{,}24915$ jest zgodne z~warunkiem
kwantowania $B_{\rm tail}(g_0^e) = 0$
(prop.~\ref{prop:J-quantization}, odchylenie $1{,}6\%$ od $g_0 = 1{,}2301$).
\end{remark}

% -------------------------------------------------------------------
\subsection{Amplitudy ogona i~profil trzech generacji}\label{app:J2-profile}
% -------------------------------------------------------------------

\begin{proposition}[Profile trzech generacji z~$\varphi$-FP]\label{prop:J2-profiles}
Z~$g_0^* = 1{,}24915$ i~zasady selekcji $g_0^{(n)} = \varphi^n \cdot g_0^*$:
\begin{center}\small
\begin{tabular}{@{}lcccc@{}}
\toprule
Generacja & $g_0$ & $A_{\rm tail}$ & $(A/A_e)^4$ & PDG \\
\midrule
$e$ ($n=0$) & $1{,}24915$ & $0{,}2988$ & $1$ & $1$ \\
$\mu$ ($n=1$) & $2{,}02117$ & $1{,}1331$ & $206{,}77$ & $206{,}768$ \\
$\tau$ ($n=2$) & $3{,}27032$ & $2{,}3698$ & $3955{,}1$ & $3477{,}48$ \\
\bottomrule
\end{tabular}
\end{center}
Warunek brzegowy $g(r \to \infty) \to 1$ spe\l{}niony:
$|g(R_{\rm max}) - 1| < 0{,}05$ dla wszystkich trzech generacji
(test T1, \texttt{ex106}).

\medskip\noindent
\textbf{Klucz}: Stosunek $r_{21}$ jest \textbf{dok\l{}adny} ($0{,}0001\%$).
Stosunek $r_{31}$ z~zasady $g_0^\tau = \varphi^2 g_0^*$ daje
odchylenie $13{,}7\%$ od PDG --- rz\k{a}d wielko\'sci poprawny,
ale predykcja nie jest precyzyjna.

\textbf{Status}: \statuslabel{Propozycja} ($r_{21}$: twierdzenie numeryczne;
$r_{31}$: hipoteza z~$\varphi^2$, problem O-J3 otwarty).
\end{proposition}

% -------------------------------------------------------------------
\subsection{Wykładnik skalowania i~definicja masy}\label{app:J2-skalowanie}
% -------------------------------------------------------------------

\begin{proposition}[Empiryczny wyk\l{}adnik skalowania]\label{prop:J2-nu}
W~zakresie $g_0 \in (g^* + 0{,}05,\; g^* + 1{,}5)$ amplituda ogona
podlega przybli\.zonemu skalowaniu:
\begin{equation}\label{eq:J2-power-law}
  A_{\rm tail}(g_0) \;\approx\; 0{,}80 \cdot (g_0 - g^*)^{1{,}36},
  \qquad R^2 = 0{,}80.
\end{equation}
Wyk\l{}adnik $\nu \approx 1{,}36$ odbiega od wcze\'sniejszego oszacowania
$\nu \approx 4{,}12$ (Dodatek~\ref{app:ogon-masy}, stw.~\ref{prop:J-scaling}):
ta r\'o\.znica wynika z~innego parametru skalowania:
\begin{itemize}[nosep]
  \item \ref{prop:J-scaling}: $A_{\rm tail} \propto g_0^{4{,}12}$
    (fit wzgl\k{e}dem $g_0$, nie $g_0 - g^*$),
  \item \eqref{eq:J2-power-law}: $A_{\rm tail} \propto (g_0 - g^*)^{1{,}36}$
    (fit wzgl\k{e}dem odleg\l{}o\'sci od bariery duchowej).
\end{itemize}
Obie parametryzacje s\k{a} numerycznie sp\'ojne; \.zadna nie jest
dok\l{}adnym power-law ($R^2 < 1$), co wskazuje, \.ze
$A_{\rm tail}(g_0)$ jest funkcj\k{a} z\l{}o\.zon\k{a}
(problem otwarty O-J1).

\textbf{Status}: \statuslabel{Propozycja} (numeryczna, \texttt{ex106}, test T5).
\end{proposition}

\begin{remark}[Sta\l{}a masy $c_M$ i~definicja formalna]\label{rem:J2-cM}
Numeryczne obliczenie masy z~pe\l{}nego Lagrangianu
$\mathcal{L} = K_{\rm sub}(g)\,g'^2/2 + V(g)$ (z~$K_{\rm sub}(g) = g^2$,
sekcja~\ref{sec:ghost-resolution}) daje:
\begin{align}
  M_{\rm full}(e) &= -0{,}306, \quad
  M_{\rm full}(\mu) = +12{,}80, \\
  c_M &\equiv M_{\rm full}(e) / A_e^4 = -38{,}4.
\end{align}
Ujemno\'s\'c $c_M$ oznacza, \.ze potencja\l{}owy wk\l{}ad $V(g) - V(1)$
dominuje nad kinetycznym $K_{\rm sub} g'^2/2$ w~rdzeniu solitonu.
\textbf{Masa fizyczna pochodzi z~kwantowego trybu zerowego}
(prop.~\ref{prop:J-zero-mode-mass}), nie z~klasycznej energii.

Stosunek $|M_{\rm full}(\mu)/M_{\rm full}(e)| \approx 42$ potwierdza,
\.ze klasyczna masa \textbf{nie reprodukuje} $r_{21} = 206{,}77$
--- masa jest efektem kwantowym proporcjonalnym do $A_{\rm tail}^4$.

\textbf{Status}: problem otwarty O-J4 (formalna definicja $c_M > 0$
wymaga regularyzacji).
\end{remark}

% -------------------------------------------------------------------
\subsection{Sp\'ojno\'s\'c z~K$_{\rm sub}$(g) = g$^2$}\label{app:J2-Ksub}
% -------------------------------------------------------------------

\begin{proposition}[Sp\'ojno\'s\'c solver\'ow: regularyzowany vs substratowy]%
\label{prop:J2-Ksub-compat}
Dla elektronowego $g_0^* = 1{,}24915$, amplitudy ogona
obliczone dwoma solverami:
\begin{itemize}[nosep]
  \item \textbf{Regularyzowany} ($f(g) = 1 + 2\alpha\ln g$,
    odbicie elastyczne przy $g^*$): $A_{\rm reg} = 0{,}2988$.
  \item \textbf{Substratowy} ($K_{\rm sub}(g) = g^2$,
    sekcja~\ref{sec:ghost-resolution}): $A_{\rm sub} = 0{,}2727$.
\end{itemize}
Odchylenie: $8{,}7\%$ --- sp\'ojne w~granicach systematyki
numerycznej (r\'o\.zne sposoby traktowania bariery).
Oba solvery daj\k{a} ten sam rz\k{a}d wielko\'sci $A_{\rm tail}$
i~identyczn\k{a} struktur\k{e} profilu.

\textbf{Status}: \statuslabel{Propozycja} (\texttt{ex106}, test T6).
\end{proposition}

% -------------------------------------------------------------------
\subsection{Zaktualizowany status problem\'ow otwartych}\label{app:J2-open}
% -------------------------------------------------------------------

\begin{center}\footnotesize
\begin{tabular}{@{}p{4.5cm}lp{5cm}@{}}
\toprule
\textbf{Problem} & \textbf{Status} & \textbf{Wynik sesji v41} \\
\midrule

O-J1: Analityczne $A_{\rm tail}(g_0)$
  & \statuslabel{Program}
  & $\nu \approx 1{,}36$ (wzgl.\ $g_0-g^*$); brak prostego power-law;
    potrzebna metoda matched asymptotic \\[3pt]

O-J2: Zasada selekcji $g_0^\mu$
  & \statuslabel{Twierdzenie}
  & \textbf{ZAMKNI\k{E}TY} przez $\varphi$-FP (def.~\ref{def:J2-phi-FP},
    twierdzenie~\ref{thm:J2-FP}): $g_0^\mu = \varphi \cdot g_0^*$ \\[3pt]

O-J3: Hierarchia $r_{31}$ (tau)
  & \statuslabel{Hipoteza}
  & $g_0^\tau = \varphi^2 g_0^*$ daje $r_{31} = 3955$
    (13{,}7\% od PDG); kierunek poprawny,
    precyzja niewystarczaj\k{a}ca \\[3pt]

O-J4: Formalna definicja $c_M > 0$
  & \statuslabel{Program}
  & $c_M = -38{,}4$ z~klasycznego Lagrangianu;
    masa z~kwantowego zero-mode \\

\bottomrule
\end{tabular}
\end{center}

\paragraph{\L{}a\'ncuch wyprowadze\'n \'Scie\.zki~9 (zaktualizowany).}
\begin{enumerate}[nosep,label=\textbf{P9.\arabic*.}]
  \item \textbf{Aksjomat N0-5}: $\beta = \gamma$ $\Rightarrow$ $V(g) = g^3/3 - g^4/4$
    (sek.~\ref{ssec:warunek-prozniowy}).
  \item \textbf{ODE solitonu}: $f(g)g'' + (2/r)g' = V'(g)$, ogon oscylacyjny
    $\delta \sim A_{\rm tail} \sin(r)/r$
    (prop.~\ref{prop:J-quantization}, stw.~\ref{prop:N0-6-from-N0-5}).
  \item \textbf{Zero-mode}: $\mathcal{E}_{\rm linear} = 0$
    $\Rightarrow$ $M \propto A_{\rm tail}^4$
    (prop.~\ref{prop:J-zero-mode-mass}).
  \item \textbf{$\varphi$-FP}: $g_0^* = 1{,}24915$ taki, \.ze
    $(A(\varphi g_0^*)/A(g_0^*))^4 = 206{,}77$
    (twierdzenie~\ref{thm:J2-FP}).
  \item \textbf{Wynik}: $r_{21} = 206{,}77$ z~odchyleniem $0{,}0001\%$ od PDG.
\end{enumerate}

\paragraph{Liczba parametr\'ow wolnych w~\L{}a\'ncuchu P9.}
\begin{itemize}[nosep]
  \item $\Phi_0$ (skala pr\'o\.zni) --- parametr TGP, okre\'sla $a_\Gamma$, $\kappa$, $G$.
  \item $\varphi$ (z\l{}ota proporcja) --- \textbf{nie jest parametrem swobodnym}:
    warto\'s\'c $g_0^*$ jest wyznaczona dynamicznie przez
    warunek~\eqref{eq:J2-FP-def}.
  \item $g_0^*$ jest rozwi\k{a}zaniem r\'ownania $R(g_0) = 0$
    --- \textbf{nie jest parametrem}, lecz wynikiem.
\end{itemize}
Efektywnie: \L{}a\'ncuch P9 ma \textbf{zero parametr\'ow wolnych}
ponad og\'olne parametry TGP ($\Phi_0$).
Stosunek mas $r_{21}$ jest \emph{predykcj\k{a}},
nie dopasowaniem.
